% Scalar price of the finite-volume Bill response and its root receipt.
\section{The paid-response price bound}
\label{sec:ex-lattice-trace-bill-prices}

We compute the quantitative root-price bound attached to a paid finite-volume
response. This is the numerical input used for the two component errors in
Section~\ref{sec:ex-lattice-trace-probe-assembly}. The computation has two
parts: a coefficient-size estimate from the paid label and its approximation
receipt, and a negative-rate bound on the native root price.

For each frame index $i$, set
\[
 J_i=\operatorname{basin}(i),\qquad K_i=\operatorname{summit}(i),
 \qquad t_i=\widehat\beta_{K_i},\qquad
 v_i=\left(\frac{A_i}{A_{J_i}}\right)^d,
 \qquad f_i=\operatorname{fly}_d(J_i),
\]
where $A_j$ is the field dilation at level $j$, and the basin, the summit
and all scales belong to one fixed selected family and its selected
schedule. Let
\[
 s_i=\operatorname{stakes}(i,J_i)
    =\exp\!\left(\Bigl(\frac{A_i}{A_{J_i}}\Bigr)^{4}
        \operatorname{Treasury}(J_i)+5\right),
 \qquad
 \gamma_i=\operatorname{gain}(K_i,J_i)
    =\exp\!\left(\frac{S_{K_i}-S_{J_i}}{8}\right)
\]
be the local stake price and the gain from level $K_i$ down to $J_i$,
where $S_j$ is the scale coordinate. The
source's native root allowance for a grade-$d$ response has the exact form
\begin{equation}
 n_{d,i}=\frac{s_i}{t_i^{-1}\left(
     \dfrac{f_i/2}{A_{J_i}^{-d}(t_i^3/\gamma_i)}\right)}.
 \label{eq:ex-trace-native-price}
\end{equation}
The norm of the inverse grade-$d$ dilation is $A_{J_i}^{-d}$, which is the
factor appearing in this expression. Simplifying gives
\begin{equation}
 A_i^d n_{d,i}
   =\frac{2s_i v_i t_i^4}{f_i\gamma_i}.
 \label{eq:ex-trace-priced-native}
\end{equation}
Indeed, the denominator in \eqref{eq:ex-trace-native-price} is
$f_i\gamma_i A_{J_i}^d/(2t_i^4)$, and $A_i^d/A_{J_i}^d=v_i$.
Throughout this calculation the scale receipts ensure
$A_i,A_{J_i},t_i,f_i,\gamma_i>0$ and $G_i\to\infty$.

The scalar rate receipts are expressed by tickets. For a positive sequence
$F_i$, a rate-zero ticket means that for every $\epsilon>0$, eventually
$0\leq F_i\leq\exp(\epsilon G_i)$, where
$G_i=\operatorname{Gate}(\widehat\beta_i)$. If two sequences have tickets
of rates $\alpha$ and $\alpha'$, their product has rate $\alpha+\alpha'$:
choose the two ticket bounds at exponent $\alpha+\eta$ and
$\alpha'+\eta$, multiply them, and then let $2\eta$ be smaller than the
requested excess over $\alpha+\alpha'$. Nonnegative sums of equal-rate
sequences retain that rate by the same argument and the elementary bound
$u+v\leq2\max(u,v)$.

The following five scale estimates are inputs to this section; they are
properties of the selected schedule, proved in the source from its
selection rules and simplicity of the gauge group, and are not derived
here. Four are rate-zero tickets:
\[
 s_i,\qquad v_i,\qquad t_i^4,\qquad f_i^{-1};
\]
the fifth is the gain estimate $\tau_i(500)\le\gamma_i$ for large $i$
stated below.
The first is the stake receipt. The second is the dilation-ratio receipt at
the selected basin. The third is the scale-at-summit receipt. For the fourth,
$J_i\leq K_i$ and the inverse-flux receipt applies to the flight scale at
$J_i$. Multiplication, including the constant factor $2$, therefore gives
a rate-zero ticket for
\[
 B_i=\frac{2s_i v_i t_i^4}{f_i}.
\]
The gain estimate supplies eventually
$\tau_i(500)\leq\gamma_i$, where
\[
 \tau_i(q)=\exp\!\left(\frac{q}{10^4}G_i\right).
\]
The ticket bound at exponent $70/10^4$ gives $B_i\leq\tau_i(70)$
eventually. Dividing by the gain and using multiplicativity yields
\begin{equation}
 A_i^d n_{d,i}=B_i/\gamma_i
   \leq\tau_i(70)/\tau_i(500)=\tau_i(-430).
 \label{eq:ex-trace-native-rate}
\end{equation}
This proves the exact negative-rate price needed below for all sufficiently
large $i$.

\begin{lemma}[Column size from a paid label]
\label{lem:ex-trace-bill-column-size}
Suppose a grade-$d$ paid label $y$ satisfies
\[
 \|y\|_{\mathrm{coeff}}\leq1,\qquad
 \|y-\operatorname{Col}(x)\|_{\mathrm{coeff}}
      \leq\vartheta_K^2\,\operatorname{Age}_K(-1/4),
\]
where the selected threshold obeys $0<\vartheta_K\leq1$ and
$\operatorname{Age}_K(-1/4)\leq1$. Then
\[
 \|\operatorname{Col}(x)\|_{\mathrm{coeff}}\leq2.
\]
\end{lemma}
\begin{proof}
The triangle inequality gives
\[
 \|\operatorname{Col}(x)\|_{\mathrm{coeff}}
   \leq\|y\|_{\mathrm{coeff}}
       +\|\operatorname{Col}(x)-y\|_{\mathrm{coeff}}
   \leq1+\vartheta_K^2\operatorname{Age}_K(-1/4)\leq2.
\]
\end{proof}

\begin{proposition}[Paid label yields the scaled component price]
\label{prop:ex-trace-bill-root-price}
Fix the volume index $N$, the period $p$ and the grade $d$, and let
$K=K_i\le N$. Let $g$ be a scale-$K$ observable, $y$ a grade-$d$
coefficient and $x$ a public grade-$d$ label, and let
$\operatorname{Col}(x)=\mathcal C_K(x)$ be the scale-$K$ response column of
$x$ defined in Section~\ref{sec:ex-lattice-trace-coffer-supplier} for the
same $N,p,d$. Suppose $(y,x)$ satisfies the hypotheses of
Lemma~\ref{lem:ex-trace-bill-column-size}. Assume its paid root receipt:
whenever the window of the selected schedule from $i$ to $J_i$ is open,
that is, $i\le J_i$ and $\sigma_{J_i}^2\le A_i/A_{J_i}$ with $\sigma_j$ the
cover size of the selected clock at level $j$,
the reference level of grade $d$ is at most $J_i$, $J_i<K_i$, the scale
$i$ has room for a radius $l$ at scale $K_i$, and $g$ is interior to the
window of radius $l$,
\[
 \mathcal R_{i,p}\!\left(
   \mathcal B\bigl(A_{K_i}^{-d}g\bigr)-\mathcal O(x)\right)
       \leq n_{d,i},
\]
where $\mathcal R_{i,p}$ is the root functional of the tiling of the
volume-$N$ lattice by scale-$i$ cells, $\mathcal B$ the extension from
scale $K_i$ to volume $N$, and $\mathcal O(x)$ the volume-$N$ observation
of the label $x$ at the origin. Assume that $i$ is so large that
\eqref{eq:ex-trace-native-rate} holds.
Then the paid column has size at most $2$. For
$q_i=(A_i/A_{K_i})^d$, linearity and homogeneity give the scaled error
bound
\[
 \mathcal R_{i,p}\!\left(
   \mathcal B(q_i g)-\mathcal O(A_i^d x)\right)
      \leq\tau_i(-430).
\]
For all sufficiently large $i$, this is the paid-component error used in the
probe assembly, uniformly over $N$, $p$, the radii $l$ and the observables
satisfying the stated window, room and interiority conditions.
\end{proposition}
\begin{proof}
The coefficient conclusion is Lemma~\ref{lem:ex-trace-bill-column-size}.
The inverse dilation satisfies $q_i g=A_i^d(A_{K_i}^{-d}g)$. By
real-linearity of the finite-volume observation and pullback,
\[
 \mathcal B(q_i g)-\mathcal O(A_i^d x)
   =A_i^d\left(\mathcal B(A_{K_i}^{-d}g)-\mathcal O(x)\right).
\]
Since $A_i>0$, homogeneity of the root seminorm and the assumed paid receipt
bound this by $A_i^d n_{d,i}$. Equation~\eqref{eq:ex-trace-native-rate}
then gives the stated $\tau_i(-430)$ bound, uniformly over all allowed
windows and observables.
\end{proof}

The label-producing receipt is obtained in the source from a funded atomic
activity and a scale-dependent resource transfer. The mathematical inputs
are not just a coefficient bound: the atomic activity carries a bounded
charge, a paid profile, and a density identity; its finite-volume conversion
also uses the price-duality estimate against every allowed test. The source
construction of those atomic and dual-price inputs remains part of the full
trace supplier. This section discharges the coefficient and rate arithmetic
once they are available.
